%!TeX root = main.tex
%!encoding = UTF-8 Unicode

\section{Switching Event Separation}
In the analysis of the switching behavior of power electronics it is common practice to analyze the four switching events separately - active-off/passive-on and active-on/passive-off \cite{fritz_toward_2020}. 
Here, the active device is the switch, which initiates the switching transient while the passive switch is turned off and the body diode operates during the commutation process.
This separated analysis is also desirable for the EMC analysis to identify the switching event that is dominating the frequency spectrum for the respective frequency band of interest.
This is enabled by-and one of the main advantages-of the measurement directly at the semiconductor ports. 
Conventional EMC measurements at LISNs or antennas can only measure the conglomerated noise of all switching events and noise sources in the system.

However, to achieve this separated analysis in the frequency domain, a method introduced in \cite{kragl_impact_2024} must be used, called snap-and-mirror function.

\begin{figure}
	\begin{subfigure}[t]{0.48\textwidth}
		\centering
		\import{./figures/}{3_Methods/SnapNMirror/DPT_TD.tex}
		\subcaption{Active and passive switching signals in time domain.}
	\end{subfigure}
	\hfill
	\begin{subfigure}[t]{0.48\textwidth}
		\centering
		\import{./figures/}{3_Methods/SnapNMirror/DPT_FD.tex}
		\subcaption{Spectrum of active and passive switch.}
	\end{subfigure}
	
	\begin{subfigure}[t]{0.48\textwidth}
		\centering
		\import{./figures/}{3_Methods/SnapNMirror/Snapped_TD.tex}
		\subcaption{Snapped and mirrored time domain signals.}
	\end{subfigure}
	\hfill
	\begin{subfigure}[t]{0.48\textwidth}
		\centering
		\import{./figures/}{3_Methods/SnapNMirror/Snapped_FD.tex}
		\subcaption{Spectrum of snapped and mirrored signals.}
	\end{subfigure}
	\label{fig:SnM_snapped_FD}
	\vspace{1cm}
\end{figure}

% \begin{figure}[h!]
% 	\centering\import{./figures/}{2_SnapNMirror/DPT_TD.tex}
% 	\caption{Principle functionality of snap and mirror function for the separated EMC analysis of all four switching events.}
% 	\label{fig:SnapNMirror}
% \end{figure}



\begin{figure}[h!]
	\centering\import{./figures/}{Snap_n_Mirror_vertical_v3_tex.pdf_tex}
	\caption{Principle functionality of snap and mirror function for the separated EMC analysis of all four switching events.}
	\label{fig:SnapNMirror}
\end{figure}

This function identifies the mid-points of each switching cycle \mi{t}{mid,n} and mirrors them around this axis, which can be formulated as

\begin{equation}
	\small 
	\mathrm{s}_\mathrm{new}(t) = 
	\begin{cases}
		s(t) 					 		& \text{, for } \mi{t}{mid,n}~~~   < t < \mi{t}{mid,n+1} \\
		s(-t + 2\cdot\mi{t}{mid,n+1}) 	& \text{, for } \mi{t}{mid,n+1} < t < \mi{t}{mid,n+2} \\
		s(t) 					 		& \text{, for } \mi{t}{mid,n+2} < t < \mi{t}{mid,n+3} \\
		s(-t + 2\cdot\mi{t}{mid,n+3}) 	& \text{, for } \mi{t}{mid,n+3} < t < \mi{t}{mid,n+4} \\
		...\\
	\end{cases}
	\label{eq:snap}
\end{equation}

This is done for both the turn-on and turn-off event for both the active and the passive switch, thereby creating four artificial signals out of the two measured voltages of active and passive switch.

Because only the signal amplitude in the frequency domain is analyzed this is a valid operation since the mirroring and the time shift of a signal has only influence on the phase of the fourier transform.

\begin{equation}
	\mathcal{F}\{s(-t)\} = S(-f) = S^\ast(f)
\end{equation}

\begin{equation}
	\mathcal{F}\{s(t-\mi{t}{0})\} = S(f)e^{-j\omega \mi{t}{0}}
\end{equation}